% Plan: development/outline.md, section 6, item E. Lemma lem:audit-display,
% Propositions prop:spring-opt and prop:emfl-enclosure and Figure
% fig:audit-margins moved here from Section 7 (development/moves/m-points-audit.md).
% Round-2 revision (adjudication G6-01, G6-06, G6-10): S4.1 holds the aggregate
% statements folded out of Section 7.5 (move M5) and S4.4 the best listed primal
% value of spring (folded out of Section 7.4, move M4); run records of the clean-copy
% reruns moved to development/moves/r2-G6.md (for artifact/RUNS.md).
\section{Audit details}\label{app:audit}

All page data of the audit are from the snapshot of 2026-09-30 (\cref{sec:audit-screen}), and ``separately written'' has the meaning of \cref{sec:semantics-protocol}.

\subsection{Data, screen and classification}\label{app:audit-screen}

\paragraph{Pages and points.}
We cut the 1,633 stored instance pages before the embedded GAMS listing and parsed them into points (displayed value, listed violation, section, date) and per-solver dual bounds (value, solver, date).
1,366 instances minimize and 264 maximize; \inst{camcns}, \inst{gancns} and \inst{korcns} have no objective and were skipped, and \inst{fct}, which has a page but no OSIL file, is in no flagged pair.
Two separately written parsers agree on all 2,816 points, 11,086 bounds and 11,431 fields of the instance list.
Every number is kept as its displayed string, which has at most eight decimals, and read as an exact rational.
The 56 flagged points are MINLPLib's \texttt{.sol} files; as MINLPLib specifies, a variable missing from such a file is 0, and the entry of the objective variable is ignored when the OSIL form has eliminated it.

\paragraph{Screen.}
No point with a listed violation above $10^{-5}$ lies below any listed dual, so the violation filter of the screen (\cref{sec:audit-screen}) removes nothing.
We evaluated each flagged point in 40-digit arithmetic from the decimal strings of the OSIL and \texttt{.sol} files.
Each of the 56 points has a violation of at most $10^{-6}$ by this evaluation or a listed violation of at most $10^{-6}$.
The evaluated violations agree with the listed ones within a factor of about ten, or both are below $3\cdot10^{-10}$, with one large exception: \inst{nd_netgen-2000-3-4-b-a-ns_7} p2 is listed with $10^{-12}$ but violates a cone row by $1.70\cdot10^{-5}$; in binary64, the value of $s^2-t^2$ with $s,t\approx1.45\cdot10^{6}$ loses all its digits.
\paragraph{Classes.}
The classes of \cref{sec:audit-screen} are decided for each pair (bound $s$, flagged point $p$) with the margin $\dval{s}-\varphi_{\mathrm{hi}}$, where $[\varphi_{\mathrm{lo}},\varphi_{\mathrm{hi}}]$ encloses the objective value of an exactly feasible point proved to exist near $p$ (\cref{app:audit-existence}).
Class (ii) means $\dval{s}\le L$ for a rigorous lower bound $L$ on $\vstar$, class (ii$'$) means $\varphi_{\mathrm{lo}}\ge\dval{s}$, and class (iii) includes every pair whose existence proof failed; a class (i) bound is \emph{gross} if its margin also exceeds $10^{-6}|\dval{s}|$ and \emph{tolerance-scale} otherwise.
An (instance, solver) pair receives the strongest class over its flagged points, in the order (i), (i-r), (ii), (ii$'$), (iii).
The audit implementation drew the line between (i) and (i-r) at the slack $\rho(s)=\max\bigl(\tfrac12\dunit{s},\tfrac12\cdot10^{\lfloor\log_{10}|\dval{s}|\rfloor-9}\bigr)$, which equals $\tfrac12\dunit{s}$ for all 158 flagged pairs.
Both rules give the same classes, because no margin of a flagged pair lies between $0.44\,\dunit{s}$ and $1.115\,\dunit{s}$: the closest decisions are a class (i) margin of $1.115\,\dunit{s}$ ($2.23\rho(s)$) and a class (i-r) margin of $0.436\,\dunit{s}$ ($0.871\rho(s)$), both computed in exact rational arithmetic.
\Cref{tab:audit-classes} gives the counts by class and solver, and \cref{fig:audit-margins} the margins of all refuted and class (i-r) bounds.
The instances of each class are listed in \cref{tab:audit-pairs} (class (i)), \cref{tab:audit-ir} (class (i-r)), \cref{app:audit-emfl} (class (ii)) and the paragraph on undecided pairs below; the 63 bounds of class (ii$'$) concern \inst{alkyl}, \inst{elf}, \inst{hybriddynamic_fixedcc}, \inst{hybriddynamic_varcc}, \inst{squfl010-040persp}, \inst{squfl015-080persp}, \inst{squfl020-040persp}, \inst{squfl020-050persp}, \inst{squfl025-040persp}, \inst{squfl030-100persp} and \inst{squfl030-150persp}.

In ten of the 19 class (i) pairs, the invalid bound agrees, to all its displayed digits, with the value of a listed point that was added on or before the date of the bound and whose value lies above the proved objective value: \inst{glider100} (COUENNE and LINDO), \inst{methanol50}, the four \inst{sssd} instances, \inst{ghg_3veh} (ANTIGONE and BARON) and \inst{nd_netgen-2000-3-4-b-a-ns_7} (Gurobi).
The pages do not record how any bound was computed, and we draw no conclusion about the cause.

\begin{table}[t]
\centering\small
\caption{Classes of the 131 flagged (instance, solver) pairs, by solver (pages of 2026-09-30). Class (i): gross if the margin exceeds $10^{-6}|d|$, tolerance-scale otherwise. ``Bounds'': per-solver bounds of the solver on all 1,633 pages; the other nine solver labels (ALPHAECP, XPRESS, PQCR and six labels with 48 bounds together) have no flagged pair.}
\label{tab:audit-classes}
\begin{tabular}{@{}lrrrrrrrr@{}}
\toprule
 & & & \multicolumn{2}{c}{(i) invalid} & (i-r) invalid & (ii) proved & (ii$'$) not & (iii) \\
\cmidrule(lr){4-5}
solver & bounds & pairs & gross & tol. & as displayed & valid & refuted & undecided \\
\midrule
ANTIGONE & 1424 & 15 & 1 & 0 & 3 & 0 & 4 & 7 \\
AOA      & 11   & 4  & 0 & 0 & 0 & 0 & 0 & 4 \\
BARON    & 1485 & 21 & 1 & 0 & 2 & 4 & 10 & 4 \\
BONMIN   & 348  & 1  & 0 & 1 & 0 & 0 & 0 & 0 \\
COUENNE  & 1451 & 15 & 1 & 0 & 3 & 0 & 9 & 2 \\
CPLEX    & 389  & 9  & 0 & 1 & 0 & 0 & 8 & 0 \\
Gurobi   & 1073 & 10 & 0 & 1 & 0 & 0 & 9 & 0 \\
LINDO    & 1564 & 35 & 8 & 5 & 2 & 4 & 11 & 5 \\
SCIP     & 1565 & 20 & 0 & 0 & 2 & 4 & 11 & 3 \\
SHOT     & 1011 & 1  & 0 & 0 & 0 & 0 & 1 & 0 \\
\midrule
pairs    &      & 131 & 11 & 8 & 12 & 12 & 63 & 25 \\
instances &     & 46 & 9 & 6 & 4 & 4 & 11 & 12 \\
\bottomrule
\end{tabular}
\end{table}

\begin{figure}[t]
\centering
\includegraphics[width=\textwidth]{fig-audit-margins}
\caption{Margins $d-\varphi$ of the refuted listed dual bounds (class (i) and \inst{rocket}) and of the class (i-r) bounds, in units of the last displayed digit of $d=\dval{s}$ (horizontal) and relative to $|d|$ (vertical). $\varphi$ is the upper end of the objective enclosure at an exactly feasible point (for \inst{spring}, the proved optimum). Pairs with the same bound and point values are drawn once (19 class (i) pairs give 15 marks, 12 class (i-r) pairs give 4). Vertical lines: the slack of half a unit and the one unit of Hypothesis~H. Horizontal lines: MINLPLib's gap tolerance $10^{-6}$ and the common relative gap tolerance $10^{-4}$. Pages of 2026-09-30.}
\label{fig:audit-margins}
\end{figure}

\paragraph{Undecided pairs.}
Class (iii) mixes three situations.
For \inst{cecil_13}, \inst{hda}, \inst{heatexch_spec2} and the four \inst{rsyn} instances, the margin at the listed point is itself below half a display unit, so these pairs could at best become (i-r); the existence proofs failed on degenerate active sets or, for \inst{hda}, on a nonpolynomial active row over fixed variables.
For \inst{ex7_3_5}, \inst{ex8_2_1b}, \inst{sepasequ_complex} and \inst{sepasequ_convent}, Newton's method did not converge, the active set was degenerate, or one row could not be verified.
For \inst{oil} (\cref{sec:audit-limits}), the active equality system has 1,141 rows of rank 1,131 after all reductions, and its rows contain $\log_{10}$, $\ln$ and the power $0.8981$, so dependent rows can be removed only by a symbolic argument, which we did not attempt.
The listed violation of its point p2 is $4\cdot10^{-12}$ (our evaluation: $7.1\cdot10^{-11}$), and Gurobi lists the value of p2 as its own bound.

\paragraph{How MINLPLib aggregates.}
These are observations from the data, not documented rules.
The instance list shows a dual bound on 1,576 pages, exactly the pages with at least three per-solver entries when infeasibility claims (entries \texttt{inf}) count as entries.
Its 1,568 numeric values (the other eight are \texttt{inf}) equal, up to display rounding, the third-best per-solver entry with \texttt{inf} counted; with finite entries only, all but four match (\inst{ball_mk4_15}, \inst{chp_shorttermplan2c}, \inst{nuclear10a}, \inst{powerflow0057r}, each with one \texttt{inf} entry).
The 589 \texttt{=bestdual=} entries of \texttt{minlplib.solu} equal, for minimization, the smaller of the third-best entry and the lowest listed point value (for maximization, the larger of the third-best entry and the highest point value); 422 agree exactly and 167 within half a displayed unit.
This is consistent with the bold marking of the three best bounds \citep{vigerske2026-minlplib-documentation-database-snapshot-2026} and with the rule of \citet{vigerske2014-minlplib-2}.
We compared these aggregate values with the lower end of each proved objective enclosure, using one exactly feasible point per instance.
On the 15 class (i) instances and on \inst{rocket100}--\inst{rocket400}, the four \inst{emfl} instances, \inst{spring} and \inst{lop97icx}, neither the dual bound in the instance list nor the \texttt{=bestdual=} entry exceeds the lower end by more than display rounding.
On \inst{eniplac} and \inst{stockcycle}, the dual bound in the instance list is a six-digit (i-r) entry that exceeds a proved objective value by at least 0.083 and 0.31, which only the unproved rounding to six digits of \cref{app:audit-spring} explains; their \texttt{.solu} values are consistent with the proved points.
Otherwise the aggregate values agree with every proved point, and no solved mark is contradicted.

\begin{table}[t]
\centering\small
\caption{SHA-256 prefixes (12 hexadecimal digits) of the OSIL files used in \cref{sec:audit}, all identical to the live files at the refresh of 2026-10-02. Full hashes are in the archive.}
\label{tab:audit-inputs}
\begin{tabular}{@{}ll@{\qquad}ll@{}}
\toprule
instance & SHA-256 & instance & SHA-256 \\
\midrule
\inst{glider100} & \texttt{9eb2662396a8} & \inst{smallinvDAXr2b200-220} & \texttt{d401b1cdf468} \\
\inst{topopt-cantilever_60x40_50} & \texttt{64c5c8510630} & \inst{watercontamination0303} & \texttt{7fa81d273eea} \\
\inst{methanol50} & \texttt{7e64e24db3e0} & \inst{eniplac} & \texttt{95bcd431040f} \\
\inst{sssd20-04persp} & \texttt{b07166acea9a} & \inst{lop97icx} & \texttt{bdb9368cdb21} \\
\inst{sssd22-08persp} & \texttt{2912c22b8d99} & \inst{spring} & \texttt{9819238d7199} \\
\inst{sssd25-04persp} & \texttt{7954821c02c3} & \inst{stockcycle} & \texttt{ee74374e3d51} \\
\inst{sssd25-08persp} & \texttt{58fa3ae9a27b} & \inst{emfl050_3_3} & \texttt{697d84fb4974} \\
\inst{ghg_3veh} & \texttt{3831bb6b455d} & \inst{emfl050_5_5} & \texttt{501dd45c154f} \\
\inst{nuclear14} & \texttt{37644467e1df} & \inst{emfl100_3_3} & \texttt{730d7ee662cd} \\
\inst{nd_netgen-2000-3-4-b-a-ns_7} & \texttt{6f9b3bd97e56} & \inst{emfl100_5_5} & \texttt{eeddebe5ff1b} \\
\inst{smallinvDAXr1b150-165} & \texttt{593ae4de8591} & \inst{rocket100} & \texttt{fa26d3cbff49} \\
\inst{smallinvDAXr2b150-165} & \texttt{234295e83620} & \inst{rocket200} & \texttt{7b1892c4e946} \\
\inst{smallinvDAXr1b200-220} & \texttt{819f3a6a9f92} & \inst{rocket400} & \texttt{28b092a0fc0a} \\
\bottomrule
\end{tabular}
\end{table}

\subsection{The display hypothesis}\label{app:audit-hyp}

\begin{lemma}[when Hypothesis~H holds]\label{lem:audit-display}
Let $b\in\R$ be a reported dual bound and $s$ the string displayed for it, with $\dval{s}\neq0$.
\begin{enumerate}
\item[(a)] Suppose $\dval{s}$ arises from $b$ by finitely many steps, each a rounding (to nearest with any tie rule, or directed) or a truncation to an integer multiple of a power of ten.
Then $|b-\dval{s}|<\tfrac{10}{9}\dunit{s}$.
If at most one step changes the value, or if the last step that changes it rounds to nearest, Hypothesis~H holds.
\item[(b)] Suppose the site stores a binary64 number $b'$ with $|b'-b|\le 2^{-53}|b|$, rounds $b'$ to nearest at ten significant digits, stores the result in binary64, and displays it rounded to nearest at eight decimals without trailing zeros.
If $|\dval{s}|\le 6\cdot10^{7}$, Hypothesis~H holds.
\end{enumerate}
\end{lemma}

\begin{proof}
(a) Let $b_0=b$ and $b_k=R_k(b_{k-1})$ for $k=1,\dots,K$, where $R_k$ rounds or truncates to an integer multiple of a power of ten $w_k$, and $b_K=\dval{s}$.
Steps with $b_k=b_{k-1}$ can be dropped.
If none remains, $b=\dval{s}$ and Hypothesis~H holds because $\dunit{s}>0$.
Otherwise renumber the remaining steps $1,\dots,k$ with units $v_1,\dots,v_k$.
For $i\ge2$ the value $b_{i-1}$ is an integer multiple of $v_{i-1}$ but not of $v_i$, since otherwise step $i$ would not change it; as both units are powers of ten, $v_i>v_{i-1}$, and hence $v_i\le10^{i-k}v_k$.
A rounding or truncation to multiples of $v$ moves a number by less than $v$, and rounding to nearest by at most $v/2$.
Hence
\[
  |b-\dval{s}|\le\sum_{i=1}^{k}|b_i-b_{i-1}|<\sum_{i=1}^{k}v_i\le v_k\sum_{m\ge0}10^{-m}=\tfrac{10}{9}v_k .
\]
Since $\dval{s}=b_k$ is a nonzero integer multiple of $v_k$, $\dunit{s}\ge v_k$, which proves the first claim.
If $k=1$, then $|b-\dval{s}|<v_1\le\dunit{s}$.
If the last step rounds to nearest, then $|b-\dval{s}|<\tfrac12v_k+\sum_{i<k}v_i\le(\tfrac12+\tfrac19)v_k<\dunit{s}$.

(b) Since $\dval{s}\ne0$, also $b'\ne0$.
Let $e=\lfloor\log_{10}|b'|\rfloor$ and $v_1=10^{e-9}$, the unit of the tenth significant digit of $b'$; then $v_1>10^{-10}|b'|$.
Let $r$ be $b'$ rounded to nearest at $v_1$, let $r'=\fl(r)$, let $v_2=10^{-8}$, and let $\dval{s}$ be $r'$ rounded to nearest at $v_2$.
Then $|b'-r|\le v_1/2$, $|r-r'|\le2^{-53}|r|$ and $|r'-\dval{s}|\le v_2/2$.

\emph{Case $v_1\ge v_2$.}
Then $r$ is an integer multiple of $v_2$.
From $|r'|\le|\dval{s}|+v_2/2<2^{26}$ and the monotonicity of rounding, $|r|<2^{26}$, so the spacing of binary64 numbers near $r$ is at most $2^{-27}$ and $|r-r'|\le2^{-28}<v_2/2$.
The multiple of $v_2$ nearest to $r'$ is therefore $r$, so $\dval{s}=r$.
This is a nonzero integer multiple of $v_1$, hence $\dunit{s}\ge v_1$.
Finally $|b-\dval{s}|\le|b-b'|+|b'-r|\le2^{-53}|b|+v_1/2<v_1$, because $2^{-52}|b|\le2^{-52}|b'|/(1-2^{-53})<10^{-10}|b'|<v_1$.

\emph{Case $v_1<v_2$.}
Then $v_1\le v_2/10$, $e\le0$, $|b'|<10$, $|r|\le10$ and $|b|<10.1$.
Summing the four errors,
\[
  |b-\dval{s}|\le2^{-53}(|b|+|r|)+\tfrac{1}{20}v_2+\tfrac12v_2<v_2\le\dunit{s},
\]
where the last inequality holds because $\dval{s}$ is a nonzero integer multiple of $v_2$.
\end{proof}

The pages are consistent with rule (b): of the 13,847 finite values they display, 13,723 are fixed points of the rule ``round to ten significant digits, print at eight decimals, remove trailing zeros'', 122 others differ only in the sign of zero, and two show more digits than the rule allows.
The rule cannot be verified from the pages alone, which is why the refutations rest on the weaker Hypothesis~H.

\begin{remark}[why (b) needs a size condition]\label{rem:audit-fac1}
The page of \inst{fac1} shows $160912612.40000001$, the eight-decimal print of the binary64 number nearest to $160912612.4$.
For this string $\dunit{s}=10^{-8}$, while a reported number that rounds to $160912612.4$ at ten significant digits may differ from $\dval{s}$ by up to $0.05$; the condition $|\dval{s}|\le6\cdot10^{7}<2^{26}$ excludes such strings.
\end{remark}

\subsection{Existence certificates}\label{app:audit-existence}

\paragraph{The Krawczyk test in midpoint–radius form.}
Both implementations evaluate the following form of \cref{thm:pts-krawczyk}, whose hypotheses they check exactly.
Absolute values and inequalities of vectors and matrices are entrywise.

\begin{lemma}[midpoint–radius Krawczyk test]\label{lem:audit-krawczyk}
Let $G:\R^n\to\R^n$ be continuously differentiable on an open set containing the box $X=\{x:|x-\tilde x|\le r\}$ with $r>0$.
Suppose $|G(\tilde x)-G_c|\le G_r$, and for every $\xi\in X$ and every row $i$, $|\nabla G_i(\xi)^{\mathsf T}-(J_c)_{i,:}|\le(J_r)_{i,:}$.
Let $R\in\R^{n\times n}$ and
\[
  \beta:=|RG_c|+|R|G_r+|I-RJ_c|\,r+|R|\,J_r\,r .
\]
If $\beta<r$, then $G$ has exactly one zero in $X$.
\end{lemma}

\begin{proof}
Apply the first part of \cref{thm:pts-krawczyk}, whose proof (\cref{app:points-krawczyk}) uses only that $G$ is continuously differentiable on an open set containing $X$, with $y=\tilde x$, $C=R$, $\mathbf v=[G_c-G_r,G_c+G_r]$, $\mathbf A=[J_c-J_r,J_c+J_r]$ and $\mathbf B=[I-RJ_c-|R|J_r,\,I-RJ_c+|R|J_r]$.
Then $G(\tilde x)\in\mathbf v$, $G'(\xi)\in\mathbf A$ on $X$, and $I-RA\in\mathbf B$ for every $A\in\mathbf A$, since $|(I-RA)-(I-RJ_c)|\le|R|J_r$.
For $v\in\mathbf v$, $M\in\mathbf B$ and $x\in X$, $|\tilde x-Rv+M(x-\tilde x)-\tilde x|\le|RG_c|+|R|G_r+(|I-RJ_c|+|R|J_r)\,r=\beta$, so $K=[\tilde x-\beta,\tilde x+\beta]$ satisfies (c), and $K\subseteq\operatorname{int}X$ because $\beta<r$.
\end{proof}

\begin{corollary}[exactly feasible point]\label{cor:audit-box}
In a Krawczyk proof (\cref{sec:audit-existence}), let $E$ be the active rows, $x_N$ the fixed values of the nonbasic variables, and $G(x_B)=g_E(x_B,x_N)-t_E$, where $t_E$ holds each active row at its side.
Suppose \cref{lem:audit-krawczyk} applies on a box $X$ of basic values, every expression of the model is defined on $X\times\{x_N\}$, and
(i) every row outside $E$ satisfies its bounds on all of $X\times\{x_N\}$, where an equality row passes only if it is constant there and equal to its side;
(ii) $X$ lies within the bounds of the basic variables;
(iii) $x_N$ satisfies its bounds and integrality requirements exactly, and every integer variable is nonbasic.
Then the zero $x_B^\ast$ yields $x^\ast=(x_B^\ast,x_N)\in\feas$, and $f(x^\ast)$ lies in every enclosure of $f$ over $X\times\{x_N\}$.
\end{corollary}

\begin{proof}
$G(x_B^\ast)=0$ puts every active row at its side, (i) covers the other rows, and (ii) and (iii) cover bounds and integrality.
\end{proof}

The implementations accept an equality row outside $E$ only if it vanishes identically as a polynomial in the basic variables or its interval enclosure is a single point equal to its side; an enclosure of positive width never passes, so no equality is accepted by accident.
Enclosures of a division or a square root whose argument interval touches 0 abort the proof.
The class (i) and \inst{rocket} files use only sums, products, negation, division, squares, integer powers, $\exp$ and $\sqrt{\cdot}$.

\paragraph{Floating-point error terms.}
Both implementations compute $R$ as an approximate inverse in binary64 and evaluate the products in $\beta$ in binary64.
Let $\uround=2^{-53}$, $\eta=2^{-1074}$ and $\gamma_m=m\uround/(1-m\uround)$.

\begin{lemma}[rounding errors of matrix products]\label{lem:audit-fp}
Let $A\in\mathbb F^{n\times m}$ and $B\in\mathbb F^{m\times p}$ be binary64 matrices with $m\uround<\tfrac12$, and let $\widehat C$ be their product computed in binary64 with round to nearest, in any order of summation, with or without fused multiply–add, and without overflow.
Then $|\widehat C-AB|\le\gamma_m|A|\,|B|+m\eta$ entrywise.
\end{lemma}

\begin{proof}
Fix an entry $c=\sum_{i=1}^m a_ib_i$.
Each rounding of a product or of a fused multiply–add returns $z(1+\varepsilon)+\delta$ with $|\varepsilon|\le\uround$ and $|\delta|\le\eta/2$, and each rounding of a sum returns $z(1+\varepsilon)$, because sums are exact in the underflow range \citep[Section~2]{rump2010-verification-methods-rigorous-results-using}.
Following each term through the computation, $\hat c=\sum_i a_ib_i\theta_i+\sum_j\delta_j\theta'_j$, where each $\theta_i$ is a product of at most $m$ factors $1+\varepsilon$, each $\theta'_j$ of at most $m-1$ such factors, and there are at most $m$ terms $\delta_j$.
Hence $|\theta_i-1|\le\gamma_m$ and $|\theta'_j|\le1+\gamma_{m-1}\le2$, which gives the bound.
Without underflow this is the classical bound; see Theorem~2.1 of \citet{rump2010-verification-methods-rigorous-results-using}.
\end{proof}

The audit implementation computes $G(\tilde x)$ and the Jacobian over $X$ by forward-mode automatic differentiation in mpmath interval arithmetic at 40 digits, starting from the exact decimals, and converts the results outward to binary64 midpoint–radius pairs.
It bounds each product in $\beta$ by \cref{lem:audit-fp}, multiplies by $(1+10^{-6})^2$ to cover the roundings of the bound computations themselves (whose relative size is below $n\uround\approx4.4\cdot10^{-13}$ for $n\le4000$), and adds $10^{-300}$ per entry for underflow.
It encloses $F$ and the Jacobian over a box whose binary64 endpoints are moved outward by one unit in the last place, and it uses as $r_i$ the smaller distance from $\tilde x_i$ to the two endpoints, so the hypotheses of \cref{lem:audit-krawczyk} hold on a box inside the enclosed one.
The separately written implementation (checker V1 below) computes the same quantities in its own interval arithmetic with rational endpoints rounded outward to 200 bits: $+,-,\times,\div$ and integer powers are exact before rounding, square roots use integer square roots, and $\exp$ and $\log$ use mpmath at 90 digits, widened by a relative $10^{-60}$.
It bounds the binary64 products by \cref{lem:audit-fp} with a factor 2 on the error terms and explicit underflow terms $k\eta$, and multiplies the result by $1+10^{-12}$; its box has exact rational endpoints.
Inspection of the code confirmed both bound computations.

\paragraph{The proof methods in detail.}
The Krawczyk proof rounds and fixes the integer variables and puts continuous variables within $10^{-7}$ (relative) of a bound on it.
The active rows are all equality rows and the inequality rows that are violated or within $10^{-7}$ of a side, held at that side; rows with vanishing gradient in the free variables move to the box check.
A basis of $|E|$ continuous variables is chosen by QR factorization with column pivoting, favouring variables away from their bounds, and the other variables keep their exact decimals.
Newton's method (residuals at 40 digits) gives the centre of a box of relative radius $10^{-12}$, $10^{-10}$ or $10^{-8}$ for the Krawczyk test, and the box check of \cref{cor:audit-box} encloses every other row, exactly in the basic variables where the row is polynomial and by intervals otherwise, and the objective.
A basic variable that ends on a bound is fixed there and the attempt repeated.
The shifted Krawczyk proof first fixes, repeatedly, every continuous variable that a linear row determines alone; two linear programs (HiGHS) then give a direction that keeps the linearized equalities, strictly improves as many active inequalities and bounds as possible and changes the objective least, and a Krawczyk proof follows at a point moved along it by a distance $t$ that grows by factors of ten from three times the violation.
The linear programs only propose the point; the Krawczyk proof proves it.

The dedicated proofs end with an exact check of every row and bound of the OSIL model, except the one for \inst{topopt-cantilever_60x40_50}, which ends with the box check of the Krawczyk proof.
For \inst{watercontamination0303}, all rows are linear and the objective is quadratic; the 14 binaries are rounded, the 1,008 bounded source variables keep their decimals, and the other 106,200 variables follow exactly from equality rows with a single unknown (largest change $1.6\cdot10^{-13}$).
For \inst{nd_netgen-2000-3-4-b-a-ns_7}, the binaries are kept; the 74 flow-conservation rows are made to hold exactly by rational Gauss--Jordan elimination on flows strictly inside their bounds (largest change $2.1\cdot10^{-13}$; the two dependent rows hold exactly), then $u:=\max(u,cf^2)$ on each arc, and $s,t$ follow from their defining rows, so each rotated cone row $cf^2+s^2-t^2\le0$ holds.
For \inst{topopt-cantilever_60x40_50}, the binaries are kept and the stresses of void elements are exactly 0 in the listed points and stay 0, so their cone rows read $0\le0$; the 3,162 (p5) or 3,056 (p4) equality rows that involve the stresses of the 1,200 solid elements are linear, and a linear Krawczyk test on a QR basis of these stresses proves a unique solution in a box (largest ratio $\beta_i/r_i$ about 0.011 at relative radius $10^{-14}$); each cost variable $c_e$ is then set to the upper end of an interval enclosure of $\sum_k w_{e,k}^2$ over the box, so every solid cone row holds and the objective $\sum_ec_e$ is an exact rational.

\paragraph{The \inst{glider100} point.}
Point p2 of \inst{glider100}, an ``other point'' with listed violation $3\cdot10^{-8}$ (exact evaluation: $1.2\cdot10^{-7}$), has objective value about $-983842.26$.
The shifted Krawczyk proof moves it by $1.75\cdot10^{-10}$ (\cref{tab:audit-second}); V1 proves a nearby point with the altitude at node 1 fixed at its bound 0.
Both describe the same trajectory: 100 steps of about 621.6\,s, an altitude that drops to 0 at node 1, rises to about 64.6\,km and ends at 900\,m, and a vertical acceleration that changes sign at every node, so that the trapezoidal averages nearly cancel.
This period-2 mode of the discretization is feasible for the model as stored; its range of about 983.8\,km compares with 1.26\,km at the listed point p1.

\paragraph{Second implementations.}
Two checkers, written separately from the audit code and from each other in separate agent sessions, re-proved the 19 class (i) pairs.
Checker V1 has its own OSIL reader with exact fractions, the interval arithmetic and Krawczyk test described above, and its own exact certificates; it covered 14 pairs and re-fetched the OSIL files and points, which were byte-identical to ours.
Checker V2 has its own reader and uses exact rational constructions only; it covered the other five pairs and recomputed every count, margin and class from the data files.
Every second proof gives an exactly feasible point whose objective value is at most the audit's upper end $\varphi$ (\cref{tab:audit-second}; for \inst{watercontamination0303} the two exact values agree in all 16 digits shown), so every margin of \cref{tab:audit-pairs} holds under either implementation.

\begin{table}[t]
\centering\footnotesize
\setlength{\tabcolsep}{3pt}
\caption{Existence proofs of the 15 class (i) conflicts. $n$: size of the square system; ratio: largest $\beta_i/r_i$ of \cref{lem:audit-krawczyk} in the audit implementation. Audit proof: method of \cref{sec:audit-existence}. Second proof: checker and method. Last column: upper end of the objective enclosure of the second proof's exactly feasible point, rounded up, or its exact value truncated as marked by ``\dots''.}
\label{tab:audit-second}
\begin{tabular}{@{}>{\raggedright\arraybackslash}p{0.3\textwidth}>{\raggedright\arraybackslash}p{0.22\textwidth}>{\raggedright\arraybackslash}p{0.2\textwidth}l@{}}
\toprule
instance (point) & audit proof & second proof & second $\varphi$ \\
\midrule
\inst{glider100} (p2) & shifted Krawczyk, $n=1208$, ratio $1.0\cdot10^{-4}$ & V1, Krawczyk, $n=1209$ & $-983842.2577881223$ \\
\inst{topopt-cantilever_60x40_50} (p5) & dedicated, linear Krawczyk, ratio $0.011$ & V1, exact Gram system and Krawczyk & $10.33547427574701$ \\
\inst{methanol50} (p4) & Krawczyk, $n=1497$, ratio $2\cdot10^{-3}$ & V1, Krawczyk & $0.0079302187993$ \\
\inst{sssd20-04persp} (p3) & Krawczyk, $n\le16$ & V1, Krawczyk; V2, exact & $347691.41048398308\dots$ \\
\inst{sssd22-08persp} (p4) & Krawczyk, $n\le16$ & V2, exact & $508713.73101118680506\dots$ \\
\inst{sssd25-04persp} (p3) & Krawczyk, $n\le16$ & V2, exact & $300176.56366586825118\dots$ \\
\inst{sssd25-08persp} (p4) & Krawczyk, $n\le16$ & V2, exact & $472093.07796988413312\dots$ \\
\inst{ghg_3veh} (p2) & Krawczyk, $n=52$, ratio $7.3\cdot10^{-5}$ & V1, Krawczyk & $7.754006050057346$ \\
\inst{nuclear14} (p3) & Krawczyk, $n=434$ & V1, Krawczyk & $-1.12968744117115$ \\
\inst{nd_netgen-2000-3-4-b-a-ns_7} (p2) & dedicated, exact & V1, exact (least $u$) & $10729657.585117233\dots$ \\
\inst{smallinvDAXr1b150-165}, \inst{smallinvDAXr2b150-165} (p2) & listed-point check & V1, V2, exact & $88.10493476$ \\
\inst{smallinvDAXr1b200-220}, \inst{smallinvDAXr2b200-220} (p2) & Krawczyk, $n=1$ & V1, V2, exact & $156.604267884$ \\
\inst{watercontamination0303} (p2) & dedicated, exact & V1, exact & $207.9850348023888\dots$ \\
\bottomrule
\end{tabular}
\end{table}

\begin{certbox}[Certificate for the class (i) refutations (\cref{tab:audit-pairs})]
\certitem{Inputs}{OSIL files of the 15 instances (SHA-256 prefixes in \cref{tab:audit-inputs}); MINLPLib's \texttt{.sol} files of the listed points; the pages of 2026-09-30.}
\certitem{Computation}{listed-point checks and the dedicated proofs for \inst{watercontamination0303} and \inst{nd_netgen-2000-3-4-b-a-ns_7}: rational arithmetic \arith{E}; Krawczyk proofs, shifted or not, and the \inst{topopt-cantilever_60x40_50} certificate: interval Jacobians in mpmath \arith{I}, Krawczyk inequality in binary64 with a priori error terms \arith{F} (\cref{lem:audit-fp}).}
\certitem{Data reading}{exact decimals or outward enclosures of them, in both implementations (confirmed by inspection of the code).}
\certitem{Implementations}{the audit implementation certifies the values $\varphi$ of \cref{tab:audit-pairs}; the second proves an exactly feasible point with objective value at most $\varphi$ for every pair (\cref{tab:audit-second}).}
\certitem{Evidence level}{\evid{hand} (\cref{prop:audit-refute}) + \evid{stored}; for \inst{topopt-cantilever_60x40_50}, \evid{hand} + \evid{rerun}: its basis is not stored, so the point is regenerated, not replayed.}
\certitem{Replay}{Tier~1; seconds per point; about 3\,min per \inst{topopt-cantilever_60x40_50} point.}
\end{certbox}

\paragraph{Trust and replay.}
Of the seven Krawczyk-based refutations (\cref{sec:audit-existence}), \inst{methanol50}, \inst{nuclear14} and \inst{topopt-cantilever_60x40_50} have rational enclosures in V1, with no transcendental function; for \inst{glider100} and \inst{ghg_3veh}, the audit needs correct outward rounding of mpmath's interval $\exp$ and $\sqrt{\cdot}$, and V1 needs mpmath's $\exp$ at 90 digits to be accurate far beyond its widening of $10^{-60}$.
Negative controls behaved as expected: moving a Krawczyk centre for \inst{methanol50} by $10^{-11}$ (relative) makes the test fail at radius $10^{-12}$ and pass at radius $10^{-11}$; tightening an inactive row of \inst{ghg_3veh} by $10^{-9}$ below its value makes the box check fail at that row; and the listed-point check rejects a point once a bound cuts off one of its values.
The \inst{topopt-cantilever_60x40_50} certificate does not store its basis, and its centre file prints values at 25 to 45 digits, so it is regenerated rather than replayed (\cref{app:repro-regen}).
With one BLAS thread, pivoted QR chooses another basis, and the regeneration proves another exactly feasible point, with objective value $10.33547432783171797\ldots$; V1's point has objective value at most $10.33547427574701$.
Both lie below the displayed $\varphi\le10.33547433$ of the archived point, so the row of \cref{tab:audit-pairs} holds for all three points.

\subsection{The class (i-r) pairs and the \texorpdfstring{\inst{spring}}{spring} optimum}\label{app:audit-spring}

The OSIL model of \inst{spring} (identical to its GAMS form) has the variables $x_1\ge0.414$, $x_2\ge0.207$, $x_3\in[\lambda,0.02]$, an integer $i_4\in[1,100]$, $x_5\ge1.1$, a free $x_6$ and binaries $b_7,\dots,b_{17}$.
Its rows are
\begin{gather*}
  x_5=x_1/x_2,\qquad x_6=\frac{4x_5-1}{4x_5-4}+\frac{0.615}{x_5},\qquad \frac{2546.47908913782\,x_6x_5}{x_2^2}\le189000,\\
  x_3=\frac{K\,x_5^3\,i_4}{x_2},\qquad 2.1x_2+1000x_3+1.05\,x_2i_4\le14,\qquad x_1+x_2\le3,\\
  x_2=\sum_{k=7}^{17}c_kb_k,\qquad \sum_{k=7}^{17}b_k=1,
\end{gather*}
with $(c_7,\dots,c_{17})=(0.207,0.225,0.244,0.263,0.283,0.307,0.331,0.362,0.394,0.4375,0.5)$, and the objective is $(a_0+a_1i_4)\,x_1x_2^2$.

\begin{proposition}[\inst{spring}]\label{prop:spring-opt}
With $a_0=1.570796327$, $a_1=0.7853981635$, $c=0.283$, $\lambda=1.78571428571429\cdot10^{-3}$, $K=6.95652173913044\cdot10^{-7}$ and $x_5^\ast=\bigl(\lambda c/(9K)\bigr)^{1/3}$, the optimal value of \inst{spring} is $\vstar=(a_0+9a_1)\,c^{3}\,x_5^\ast=0.846245665643154281251664635037\ldots$, attained with $i_4=9$, $x_2=0.283$ and $x_3=\lambda$.
The five displayed bounds 0.84624567 equal $\vstar$ rounded to eight decimals.
\end{proposition}

\begin{proof}
The last two rows and integrality give $x_2=c$ for exactly one $c\in\{c_7,\dots,c_{17}\}$.
Fix $i_4=n$ and $x_2=c$.
The equality rows then determine $x_1=cx_5$, $x_3=Knx_5^3/c$ and $x_6$ as functions of $x_5$, and the objective becomes $(a_0+a_1n)c^3x_5$, which is increasing in $x_5$.
Each remaining condition is a closed interval in $x_5$: the bounds give $x_5\ge1.1$ and $x_5\ge0.414/c$; the bounds on $x_3$ give $(\lambda c/(Kn))^{1/3}\le x_5\le(0.02c/(Kn))^{1/3}$; the fifth row gives $x_5^3\le(14-2.1c-1.05cn)\,c/(1000Kn)$; and the sixth gives $x_5\le(3-c)/c$.
For the third row, $x_6x_5=h(x_5)+0.615$ with
\[
  h(x)=\frac{x(4x-1)}{4(x-1)}=(x-1)+\frac74+\frac{3}{4(x-1)},\qquad h''(x)=\frac{3}{2(x-1)^3}>0\quad(x>1),
\]
so the row reads $h(x_5)\le189000c^2/2546.47908913782-0.615$, and the sublevel set of the convex function $h$ is an interval.
Hence the feasible values of $x_5$ for $(n,c)$ form a closed interval $I(n,c)$, possibly empty, and
\[
  \vstar=\min_{n,c}\,(a_0+a_1n)\,c^3\min I(n,c),
\]
the minimum over the nonempty intervals.

At $(n,c)=(9,0.283)$ the condition $x_3\ge\lambda$ gives $x_5\ge x_5^\ast=(\lambda c/(9K))^{1/3}\approx4.3217$, which exceeds the other lower bounds $1.1$ and $0.414/0.283\approx1.463$.
At $x_5^\ast$ every other condition holds: the third row evaluates to about $187991.19\le189000$, the fifth to about $5.054\le14$, the sixth to about $1.506\le3$, and $x_3=\lambda\le0.02$.
So $x_5^\ast\in I(9,0.283)$ and every element of $I(9,0.283)$ is at least $x_5^\ast$; hence $\min I(9,0.283)=x_5^\ast$, with value $(a_0+9a_1)\,0.283^3\,x_5^\ast$, attained by $i_4=9$, $b_{11}=1$, $x_5=x_5^\ast$ and the other variables defined by the equality rows.

The comparison with the other $1{,}099$ assignments is computer-assisted and exact.
One program enumerates all 1,100 assignments, brackets every cube root between rationals checked by exact cubing, and finds 890 assignments infeasible and 210 feasible, with every other minimum above the value at $(9,0.283)$; the second smallest is $0.8592755352970280\ldots$, at $(5,0.307)$.
A separately written program excludes 723 assignments by the bounds on $x_3$ and the fifth and sixth rows, and shows with exact comparisons (and rational brackets of cube roots and of $\sqrt3$) that no assignment reaches an objective value at or below a rational $T$ slightly below the value at $(9,0.283)$.
Together with an exactly feasible point, this gives, without the first program,
\begin{multline*}
  \vstar\in[0.84624566564315428125166463503714129527532857794010,\\
  0.84624566564315428125166463503714129527534815925597].
\end{multline*}
This enclosure rounds to 0.84624567 at eight decimals.
\end{proof}

\begin{certbox}[Certificate for \cref{prop:spring-opt}]
\certitem{Inputs}{\inst{spring} OSIL file (SHA-256 \texttt{9819238d7199}\dots).}
\certitem{Computation}{reduction to one variable per assignment (a written argument, checked exactly against the model at 200 random rational points); exact enumeration of the 1,100 assignments with cube roots bracketed by rationals \arith{E}.}
\certitem{Data reading}{exact decimals.}
\certitem{Implementations}{two separately written exact programs, one enumerating all minima and one excluding assignments by bounds; a 60-digit scan reproduces the 210 feasible assignments and the three smallest minima (evidence only).}
\certitem{Evidence level}{\evid{hand} + \evid{stored}.}
\certitem{Replay}{Tier~1; seconds.}
\end{certbox}

The best listed primal value of \inst{spring}, 0.8462441, lies more than $1.5\cdot10^{-6}$ below $\vstar$.

\paragraph{The other class (i-r) pairs.}
\Cref{tab:audit-ir} lists all 12 pairs.
For \inst{lop97icx} and \inst{stockcycle} the listed point is exactly feasible as listed.
For \inst{eniplac}, the listed point violates 31 equality rows by $4.8\cdot10^{-14}$ to $2.5\cdot10^{-10}$; an exactly feasible point keeps 36 listed decimals and the binaries and recomputes the other 81 variables from equality rows that are linear in them.
All these proofs are in exact rational arithmetic, by a checker written separately from the audit code, and agree with the audit's enclosures.
The last column of \cref{tab:audit-ir} shows that the page's own display rule (\cref{lem:audit-display}(b)) cannot explain the conflicts on \inst{eniplac}, \inst{lop97icx} and \inst{stockcycle}: a stored value at or below the proved objective would display with more digits.
Rounding to six significant digits before storage would explain them: all seven displayed values equal the best listed point value of their instance rounded half up to six significant digits, and the same pages show other bounds with nine or ten digits (for example BARON's $-132117.083$ on \inst{eniplac}).
Among the bounds dated 2013-09-17, 59 show exactly six significant digits and only 2 show seven, while among all other bounds seven-digit values are more common than six-digit ones (471 against 263).
This supports the explanation but does not prove it.

\begin{table}[t]
\centering\footnotesize
\setlength{\tabcolsep}{3pt}
\caption{The 12 class (i-r) pairs, all dated 2013-09-17. Solvers: \inst{eniplac} COUENNE, LINDO and SCIP; \inst{lop97icx} ANTIGONE; \inst{spring} ANTIGONE, BARON, COUENNE, LINDO and SCIP; \inst{stockcycle} ANTIGONE, BARON and COUENNE. $\varphi$: objective value of an exactly feasible point (exact; ``\dots'' marks truncation). Slack: half a unit in the last displayed digit. Last column: $d-\varphi$ divided by half a unit in the tenth significant digit or the eighth decimal, whichever is coarser.}
\label{tab:audit-ir}
\begin{tabular}{@{}llllll@{}}
\toprule
instance & $d$ & $\varphi$ & $d-\varphi$ & $(d-\varphi)/\text{slack}$ & $(d-\varphi)/\text{page half-unit}$ \\
\midrule
\inst{eniplac} & $-132117.$ & $-132117.0830199887\dots$ & $0.08301\dots$ & $0.166\dots$ & $1660\dots$ \\
\inst{lop97icx} & $4099.06$ & $4099.059953600000099$ & $4.6399999901\cdot10^{-5}$ & $0.00927\dots$ & $92.7\dots$ \\
\inst{spring} & $0.84624567$ & $0.8462456656431542\dots$ & $4.3568\dots\cdot10^{-9}$ & $0.871\dots$ & $0.871\dots$ \\
\inst{stockcycle} & $119949.$ & $71969213/600$ & $187/600$ & $0.623\dots$ & $6233\dots$ \\
\bottomrule
\end{tabular}
\end{table}

\subsection{The \texorpdfstring{\inst{emfl}}{emfl} instances}\label{app:audit-emfl}

The certificate code checks, from each OSIL file, that the instance has the form
\begin{equation}\label{eq:audit-emfl}
  \min\ \sum_{k}c_kt_k\quad\text{s.t.}\quad t_k^2-\|w_k\|^2\ge0,\ t_k\ge0,\quad w_k=A_kz+e_k,\quad z_j\ge0\ (j\in P),\ z_j\ \text{free}\ (j\notin P),
\end{equation}
with rational data and $c_k\ge0$, where each component of each $w_k$ is a variable defined by one linear equality row, and that there are no other rows, no objective constant and no integer variables.
The four instances have 531, 1,875, 981 and 3,125 cones and 18, 50, 18 and 50 base variables $z$.

\begin{lemma}[weak duality for \eqref{eq:audit-emfl}]\label{lem:audit-socp}
\begin{enumerate}
\item[(a)] If a rational $\bar z$ satisfies the sign constraints and rationals $\tau_k\ge0$ satisfy $\tau_k^2\ge\|A_k\bar z+e_k\|^2$, then $(\bar z,(A_k\bar z+e_k)_k,\tau)$ is exactly feasible and $\vstar\le\sum_kc_k\tau_k$.
\item[(b)] If rational vectors $y_k$ satisfy $\|y_k\|^2\le c_k^2$, and $g:=\sum_kA_k^{\mathsf T}y_k$ satisfies $g_j\ge0$ for $j\in P$ and $g_j=0$ for $j\notin P$, then $\vstar\ge\sum_ke_k^{\mathsf T}y_k$.
\end{enumerate}
\end{lemma}

\begin{proof}
(a) holds by construction.
For (b), let $(z,w,t)$ be feasible.
From $t_k\ge0$ and $t_k^2\ge\|w_k\|^2$ we get $t_k\ge\|w_k\|$, and by the Cauchy--Schwarz inequality $c_kt_k\ge c_k\|w_k\|\ge\|y_k\|\,\|w_k\|\ge y_k^{\mathsf T}w_k$.
Summing, $\sum_kc_kt_k\ge\sum_ky_k^{\mathsf T}(A_kz+e_k)=g^{\mathsf T}z+\sum_ke_k^{\mathsf T}y_k\ge\sum_ke_k^{\mathsf T}y_k$, since $g_jz_j\ge0$ for every $j$.
\end{proof}

\begin{proposition}[\inst{emfl}]\label{prop:emfl-enclosure}
For \inst{emfl050_3_3}, \inst{emfl050_5_5}, \inst{emfl100_3_3} and \inst{emfl100_5_5}, the optimal values lie in the first-line intervals of \cref{tab:audit-emfl}, each of width below $3.4\cdot10^{-11}$.
Every listed dual bound of these instances is valid.
The best listed primal values lie below $\vstar$ by at least $1.41\cdot10^{-5}$, $1.98\cdot10^{-3}$, $2.92\cdot10^{-4}$ and $6.86\cdot10^{-6}$, so no exactly feasible point attains them; for \inst{emfl050_3_3}, which is marked solved, this is at least $1.36\cdot10^{-6}$ relative.
\end{proposition}

\begin{proof}
Apply \cref{lem:audit-socp}.
The points $\bar z$ are rational roundings of a numerical solution of the cone program; $w$ is computed exactly, and $\tau_k$ is a rational upper bound of $\|w_k\|$ obtained from an integer square root.
The vectors $y_k$ start from a numerical dual solution and are made exactly admissible: checker V2 scales each cone to a radius slightly below $c_k$, corrects negative $g_j$ by moves within single cones and applies a final scaling; the audit makes $g=0$ exactly by a rational least-norm correction and then scales all $y_k$ by one rational factor.
Both conditions of (b) and both objective values are evaluated in rational arithmetic.
This gives the enclosures of \cref{tab:audit-emfl}.
Every listed dual bound is at most the lower end, so it is valid.
The best listed primal value plus half a unit in its last displayed digit is below the lower end by the stated amounts, so no exactly feasible point attains it.
\end{proof}

\begin{table}[t]
\centering\footnotesize
\setlength{\tabcolsep}{3pt}
\caption{Enclosures of the optimal values of the \inst{emfl} instances (outward displays; first line: checker V2, second line: audit implementation). Best listed dual and best listed primal value (point, listed violation) from the pages of 2026-09-30. Shortfall: lower end minus the best listed primal value widened by half a display unit, rounded down. $^{\checkmark}$: marked solved on MINLPLib.}
\label{tab:audit-emfl}
\begin{tabular}{@{}ll>{\raggedright\arraybackslash}p{0.16\textwidth}>{\raggedright\arraybackslash}p{0.355\textwidth}l@{}}
\toprule
instance & best listed dual & best listed primal & $\vstar\in$ & shortfall $\ge$ \\
\midrule
\inst{emfl050_3_3}$^{\checkmark}$ & $10.40173999$ & $10.40173793$ (p2, $2\cdot10^{-10}$) & $[10.4017521318429,\,10.4017521318448]$ \newline $[10.4017521316,\,10.4017521319]$ & $1.41\cdot10^{-5}$ \\
\inst{emfl050_5_5} & $18.91340776$ & $18.91165289$ (p6, $10^{-8}$) & $[18.9136329557291,\,18.9136329557624]$ \newline $[18.9136329529,\,18.9136329563]$ & $1.98\cdot10^{-3}$ \\
\inst{emfl100_3_3} & $18.13262446$ & $18.13236088$ (p4, $10^{-8}$) & $[18.1326531242336,\,18.1326531242359]$ \newline $[18.1326531194,\,18.1326531244]$ & $2.92\cdot10^{-4}$ \\
\inst{emfl100_5_5}$^{\checkmark}$ & $32.63818348$ & $32.63818348$ (p1, $10^{-10}$) & $[32.6381903545115,\,32.6381903545301]$ \newline $[32.6381903513,\,32.6381903548]$ & $6.86\cdot10^{-6}$ \\
\bottomrule
\end{tabular}
\end{table}

All listed points of the four instances except p3 of \inst{emfl050_3_3} lie below the exact optimum, the farthest by about $8.1\cdot10^{-3}$ (about $4.3\cdot10^{-4}$ relative; an ``other'' point of \inst{emfl050_5_5}); points with listed violations of $4\cdot10^{-11}$ to $8\cdot10^{-10}$ are among them, so even very small violations at the apex of a cone lower the objective visibly.
Checker V1, a third implementation, encloses the optimum of \inst{emfl050_3_3} in $[10.40175213184103,\,10.40175213184476]$.
The best listed dual bound of this instance, 10.40173999, lies at least $1.213\cdot10^{-5}$ (at least $1.166\cdot10^{-6}$ relative) below the exact optimum; its solved mark is consistent with the listed primal value 10.40173793, and since the mark is defined with tolerance-feasible points, we do not call it wrong.

\begin{certbox}[Certificate for \cref{prop:emfl-enclosure}]
\certitem{Inputs}{OSIL files of the four \inst{emfl} instances (SHA-256 prefixes in \cref{tab:audit-inputs}).}
\certitem{Computation}{rational primal points and rational dual vectors checked in exact arithmetic, square roots by integer square root \arith{E}; the numerical cone solver only proposes candidates.}
\certitem{Data reading}{exact decimals.}
\certitem{Implementations}{two for all four instances (the audit implementation and checker V2), a third (checker V1) for \inst{emfl050_3_3}; the displayed enclosures, and hence their widths, rest on checker V2 alone; the wider enclosures of the audit implementation (and of checker V1) also prove the validity of every listed dual bound and the stated shortfalls, so these verdicts are verified by a separately written implementation.}
\certitem{Evidence level}{\evid{hand} (\cref{lem:audit-socp}) + \evid{rerun}: the rational dual vectors are not stored; a rerun solves the cone programs again and checks the new vectors exactly, and may give slightly different, equally valid endpoints.}
\certitem{Replay}{Tier~1; seconds to minutes per instance.}
\end{certbox}

\subsection{The \texorpdfstring{\inst{rocket}}{rocket} instances}\label{app:audit-rocket}

The \inst{rocket} instances discretize the Goddard rocket problem with $N=100$, 200 and 400 steps.
They have $6N+7$ variables (one step length, and velocity $v$, height $h$, gravity $g$, mass $m$, thrust $T$ and drag $D$ at $N+1$ nodes) and $5N+2$ equality rows.
The objective is $\min -h_N$; the drag rows are $D_i=310v_i^2\exp(500(1-h_i))$, the gravity rows $g_i=h_i^{-2}$, the thrust lies in $[0,3.5]$ and the mass in $[0.6,1]$, with $m_0=1$ and $m_N=0.6$.

The points come from a local solver.
Both proofs fix the thrust at the solver's values, compute the step and the masses exactly (the mass rows are linear once the thrust is fixed), and run a Krawczyk test on the $4N$ unknowns $(v,h,g,D)$ at radius $10^{-12}\max(1,|x|)$, with largest ratio about $1.1\cdot10^{-4}$.
The thrust is at a bound at 76, 151 and 315 of the $N+1$ nodes, and the mass equals 0.6 exactly at 64, 127 and 254 nodes; these and all other bounds are checked exactly.
The first proof uses mpmath interval arithmetic, the second the interval arithmetic and Krawczyk code of checker V1 with its own setup (\cref{tab:audit-rocket}).
The other listed bounds of these instances are consistent with the enclosures; for example, COUENNE's $-1.01283202$ on \inst{rocket100} lies about $1.3\cdot10^{-8}$ below the enclosure.
The GAMS and OSIL forms agree exactly, with $\exp$ treated as an uninterpreted function of exactly compared arguments, and archived copies show the models unchanged since 2001.
A separate agent session reran the second proof from clean copies of the models, points and code and reproduced every inclusion test, all three brackets and the numbers of pinned nodes.

\begin{table}[t]
\centering\footnotesize
\setlength{\tabcolsep}{3pt}
\caption{LINDO's bounds on the \inst{rocket} instances against exactly feasible points. Enclosures of the first proof, rounded outward; the second proof gives the same intervals except for \inst{rocket400}, where it gives $[-1.0128365294843,\,-1.0128365294821]$. Margins from the upper ends of the first proof, absolute, relative to $|d|$ and in display units, rounded down.}
\label{tab:audit-rocket}
\begin{tabular}{@{}llllll@{}}
\toprule
instance & LINDO $d$ (date) & objective enclosure & $d-\varphi\ge$ & rel.\ $\ge$ & units $\ge$ \\
\midrule
\inst{rocket100} & $-1.0128319$ (2018-05-23) & $[-1.0128320069151,\,-1.0128320069130]$ & $1.06\cdot10^{-7}$ & $1.05\cdot10^{-7}$ & $1.06$ \\
\inst{rocket200} & $-1.01283563$ (2015-03-08) & $[-1.0128356770698,\,-1.0128356770677]$ & $4.70\cdot10^{-8}$ & $4.64\cdot10^{-8}$ & $4.70$ \\
\inst{rocket400} & $-1.01283634$ (2017-09-13) & $[-1.0128365294842,\,-1.0128365294822]$ & $1.89\cdot10^{-7}$ & $1.87\cdot10^{-7}$ & $18.9$ \\
\bottomrule
\end{tabular}
\end{table}

\begin{certbox}[Certificate for the \inst{rocket} refutations (\cref{tab:audit-rocket})]
\certitem{Inputs}{OSIL files of \inst{rocket100}, \inst{rocket200}, \inst{rocket400} (\cref{tab:audit-inputs}); thrust values of the local-solver points.}
\certitem{Computation}{exact step and masses \arith{E}; Krawczyk test on $4N$ unknowns with $\exp$ enclosed by mpmath interval arithmetic (first proof, \arith{I}) or by mpmath at 90 digits with a relative widening of $10^{-60}$ (second proof); binary64 matrix products with the bounds of \cref{lem:audit-fp} \arith{F}.}
\certitem{Data reading}{exact decimals or outward enclosures.}
\certitem{Implementations}{two; the margins use the first proof's upper ends.}
\certitem{Evidence level}{\evid{hand} (\cref{prop:audit-refute}, \cref{lem:audit-display}) + \evid{stored}.}
\certitem{Replay}{Tier~1; about one second per instance.}
\end{certbox}

\subsection{Model history and model forms}\label{app:audit-history}

\Cref{tab:audit-history} summarizes the evidence that the refuted bounds were computed on today's models.
The sources are the GAMS World archive of GLOBALLib and MINLPLib~1 files (2001--2010), the MINLPLib~2 copies added to MINLPLib.jl on 2017-11-23, archived instance pages from 2018 on (each embeds the complete \texttt{.gms} model), archived statistics pages from 2014-12-09 on, and the file dates of the current files.
Old files were converted with GAMS and compared with today's model.
The comparison is exact for the GAMS text and for variables, types and bounds; for rows written differently it uses 50-digit evaluation at random points, which is evidence, not proof.
Every complete archived listing in \cref{tab:audit-history} is identical to today's \texttt{.gms} text apart from leading and trailing newlines, and the pages captured in 2018 already list the flagged bounds of 2013 to 2018 with today's values and dates.

\begin{table}[t]
\centering\footnotesize
\setlength{\tabcolsep}{3pt}
\caption{Model history of the instances with refuted bounds. ``Listing'': the complete model embedded in an archived instance page.}
\label{tab:audit-history}
\begin{tabular}{@{}>{\raggedright\arraybackslash}p{0.25\textwidth}>{\raggedright\arraybackslash}p{0.1\textwidth}>{\raggedright\arraybackslash}p{0.27\textwidth}>{\raggedright\arraybackslash}p{0.3\textwidth}@{}}
\toprule
instance & bound dates & copies before the bound & verdict \\
\midrule
\inst{glider100} & 2014, 2015 & GLOBALLib 2004 & unchanged \\
\inst{topopt-cantilever_60x40_50} & 2022 & listings 2018-09-22, 2020-01-11 (complete) & unchanged; the 2022-05-22 listing is cut at 1\,MiB and its 23,479 complete lines match; file dates 2019 and 2020 \\
\inst{methanol50}, \inst{nuclear14} & 2022 & GLOBALLib or MINLPLib~1 2001; listings 2018, 2020 & unchanged \\
\inst{nd_netgen-2000-3-4-b-a-ns_7} & 2022 & listings 2018, 2020 & unchanged \\
\inst{ghg_3veh} & 2013, 2014 & MINLPLib~1 2010; MINLPLib.jl 2017 & three products of constants later written as rounded decimals ($\le1.84\cdot10^{-15}$ relative); refutation re-proved on the old text \\
four \inst{sssd}, \inst{watercontamination0303}, four \inst{smallinvDAX} & 2014-02/03 & none (added days before the bounds) & statistics 2014-12-09 and full copies 2017-11-23 equal today's; March to December 2014 not covered \\
\inst{rocket100}, \inst{rocket200}, \inst{rocket400} & 2015--2018 & GLOBALLib 2001; MINLPLib.jl 2017 & unchanged \\
\inst{eniplac}, \inst{lop97icx}, \inst{spring}, \inst{stockcycle} & 2013 & MINLPLib~1 2001 & unchanged \\
\bottomrule
\end{tabular}
\end{table}

For \inst{ghg_3veh}, the 2010 and 2017 files write products of constants in factored form, for example $150000\cdot0.0181052631578947$; today's file writes three such products as the rounded decimals 2715.7894736842, 376.046780997472 and 28.341428570246, in six places of four rows.
The text changed between 2017-11-23 and 2018-09-22, after both flagged bounds.
The audit implementation, run on both old texts, gives an objective enclosure $[7.754006050048,\,7.754006050065]$, and a separate outward-rounded interval proof on the 2010 text gives a point with objective value $7.754006050100459$; both are far below the bounds $7.7543245$.
For the \inst{smallinvDAX} instances, older GAMS versions used a default upper bound of 100 for integer variables \citep{vigerske2026-minlplib-a-library-of-mixed}; the proving points use lots of at most 47 and 63, so the verdicts hold under either default.

\begin{corollary}[tolerances and model forms]\label{cor:audit-tolerance}
A refutation as in \cref{prop:audit-refute}(a) holds (a) for every feasibility tolerance of the solver or of the library, and (b) for the GAMS form of the instance whenever that form defines exactly the same model as the OSIL file.
\end{corollary}

\begin{proof}
Every tolerance-relaxed feasible set contains $\feas$, so its infimum is at most $f(x^\ast)<b$ (\cref{prop:sem-refutation}); two forms of the same model have the same feasible points and objective.
\end{proof}

The exact comparison of today's GAMS and OSIL forms (\cref{cor:audit-tolerance}(b)) reads both files with every decimal as a rational, expands each row and the objective into an exact rational function, with $\exp$ and $\sqrt{\cdot}$ as uninterpreted functions whose arguments are compared exactly, and compares variable names, types and bounds.
It found identical models for 14 of the 15 class (i) instances, the three \inst{rocket} instances, the four \inst{emfl} instances, \inst{spring} and \inst{eniplac}.
Changing one coefficient of \inst{sssd20-04persp}, the upper bound of $x_3$ in \inst{spring} (by $10^{-13}$ absolute), or one argument of $\exp$ in \inst{ghg_3veh} is detected.
The comparison was written and run by one implementation; on 2026-10-04 a separate agent session reran it from a clean copy and obtained the same 23 identical models, the same \inst{methanol50} exception and the same three detections.
For \inst{methanol50}, the OSIL file writes 360 objective coefficients (the constant, 104 linear and 255 quadratic ones) as binary64 products, with relative changes of at most $2.39\cdot10^{-16}$ (exact maximum $1/4196280000000000$); all 1,497 rows agree.
On the box $p_4\pm1$, which contains the proved point, the two objectives differ by at most $3.2\cdot10^{-15}$ (exact bound).
For \inst{lop97icx}, 30 objective coefficients differ in the same way; its listed point is exactly feasible in both forms, with objective $4099.0599536$ in the GAMS form, so its (i-r) verdict holds in both.
\inst{stockcycle} has no exact comparison, because its objective row is a rational function that the comparison code does not expand; GAMS conversion and evaluation at sample points agree.

\paragraph{Binary64 data.}
\Cref{cor:audit-tolerance} concerns the decimal reading of the files.
To see whether the refutations depend on it, we reran the audit's point certificates with every datum that is not a binary64 number replaced by its binary64 value (\reading{c}).
Every certificate succeeded again.
The objective upper ends changed by at most $2.13\cdot10^{-10}$ (\inst{nd_netgen-2000-3-4-b-a-ns_7}), $2.85\cdot10^{-11}$ (\inst{topopt-cantilever_60x40_50}, p5, whose certificate is regenerated) and $1.04\cdot10^{-12}$ (\inst{watercontamination0303}), and not in the first twelve digits elsewhere; the smallest margin is again $1.115\,\dunit{s}$ (\inst{smallinvDAXr1b200-220}), and the \inst{rocket} margins are again 1.06, 4.70 and 18.9 display units.
This was computed with one implementation and is not a claim of the paper.
